\begin{problem}[Warmup]{(a)}
    Suppose that $f = (x - \alpha_1)\cdots(x - \alpha_5)$. Calculate the coefficient on the $x^3$ term that results when we multiply $f$ out.
\end{problem}
\unsolved

\begin{problem}[Warmup]{(b)}
    Suppose that we have fields $F \subseteq E \subseteq K$. Which of the following groups is a subgroup of which? \[\mathrm{Gal}(K \,|\, E) \quad \mathrm{Gal}(K \,|\, F) \quad \mathrm{Gal}(E \,|\, F)\]
\end{problem}
\unsolved

\begin{problem}[Warmup]{(c)}
    What further relationship holds among the groups in Exercise b, if $K$ is a finite normal extension of $F$, and $E$ is a normal extension of $F$?
\end{problem}
\unsolved

\begin{problem}[Warmup]{(d)}
    Suppose that we have the following containments among groups: $H_1 \subseteq H_2 \subseteq \mathrm{Gal}(K \,|\, F)$. What containment relationship always holds for the fixed fields $\mathrm{Fix}(H_1)$ and $\mathrm{Fix}(H_2)$?
\end{problem}
\unsolved

\begin{problem}[Warmup]{(e)}
    Suppose that $K$ is a finite normal extension of the field $F$, and $|\mathrm{Gal}(K \,|\, F)| = p$, where $p$ is a positive prime integer. What can you say about intermediate fields $E$, where $F \subset E \subset K$?
\end{problem}
\unsolved

\begin{problem}{1}
    Consider the splitting field $K$ of $x^5 - 3$ over $\mathbb{Q}(\zeta)$, where $\zeta$ is
    the primitive fifth root of unity. Compute the Galois group
    $\mathrm{Gal}(K \,|\, \mathbb{Q}(\zeta))$. Then illustrate the correspondence between
    subfields and subgroups given by the Fundamental Theorem of Galois Theory, by drawing a
    diagram illustrating all of the subgroups of the Galois group of this splitting field
    and the corresponding order-reversed picture of all the fields between
    $\mathbb{Q}(\zeta)$ and the splitting field.
\end{problem}
\unsolved

\begin{problem}{2}
    Consider the cyclotomic polynomial
    $\Phi_5 = x^4 + x^3 + x^2 + x + 1 \in \mathbb{Q}[x]$. Then we know that the splitting
    field for $\Phi_5$ is just $\mathbb{Q}(\zeta)$, where $\zeta$ is the primitive fifth
    root of unity, and the Galois group of this splitting field is a cyclic group of order
    4 (see Theorem 47.3 and Example 47.10). Draw a diagram illustrating all of the
    subgroups of the Galois group of this splitting field and the corresponding
    order-reversed picture of all the fields between $\mathbb{Q}$ and the splitting field.
\end{problem}
\unsolved

\begin{problem}{3}
    Consider again the splitting field of $x^4 - 2 \in \mathbb{Q}[x]$, as in Exercise 45.6
    and Example 47.9. As in the previous exercise, draw a diagram illustrating all of the
    subgroups of the Galois group of this splitting field and the corresponding
    order-reversed picture of all the fields between $\mathbb{Q}$ and the splitting field.
    You should explicitly verify that each intermediate field is the fixed field of the
    appropriate subgroup.
\end{problem}
\unsolved

\begin{problem}{4}
    Which of the subgroups of the Galois group in the previous exercise are normal in the
    Galois group of the splitting field? What property do the corresponding fields have?
\end{problem}
\unsolved

\begin{problem}{5}
    Prove Theorem 48.1.
\end{problem}
\unsolved

\begin{problem}{6}
    In this exercise you check a computation from Example 48.7. In that example we showed
    that $\zeta + \zeta^6$ is a root of the polynomial $x^3 + x^2 - 2x - 1$, where $\zeta$
    is the primitive seventh root of unity. Show that $\zeta^2 + \zeta^5$ and
    $\zeta^3 + \zeta^4$ are the other two roots.
\end{problem}
\unsolved

\begin{problem}{7}
    In this exercise you provide an inductive proof for the formula for the symmetric
    polynomials, as discussed in Section 48.2. Suppose that
    $f = (x - \alpha_1)(x - \alpha_2)\cdots(x - \alpha_n) \in F[x]$, where $F$ is a field.
    Let ${(-1)}^{n - k}a_{n - k}$ be the coefficient on $x^k$ when we multiply $f$ out
    using the distributive law. Use induction on $n$ to prove that
    \[a_{n - k} = \Sigma\{\alpha_{j_1}\alpha_{j_2}\cdots\alpha_{j_k}\},\]
    where the sum is over all possible choices of $k$ distinct $\alpha_i$.
\end{problem}
\unsolved

\begin{problem}{8}
    Suppose that $f \in F[x]$ is irreducible polynomial of degree $n$ over the field $F$,
    where $F$ is a field of characteristic zero. Let $K$ be the splitting field of $f$ over
    $F$. Then we can think of $G = \mathrm{Gal}(K \,|\, F)$ as a subgroup of the
    permutation group $S_n$. Prove that $G$ is a transitive subgroup of this permutation
    group (see Exercise 34.13 for a definition).
\end{problem}
\unsolved
